\chapter*{旧目录}

\section*{第零部分\quad 绪论} 研究代数几何的理由，“子集”问题；几何的不同范畴，交换代数的必要性，部分定义函数；作者特点.先修知识，与其他课程的关系，书目列表.

\section*{第一部分\quad 平面曲线的游戏}

\noindent
\textbf{1. 平面二次曲线}\;  对 \({\mathbb{P}}^{2}\) 和齐次坐标的基本熟悉， \({\mathbb{A}}^{2}\) 与 \({\mathbb{P}}^{2}\) 的关系；参数化，每个光滑二次曲线 \(C \subset  {\mathbb{P}}^{2}\) 都是 \(\cong  {\mathbb{P}}^{1}\) .贝祖定理的简单情况:直线 \(\cap\) 与 \(d = d\) 次曲线的交点，二次曲线 \(\cap\) 与 \(d = {2d}\) 次曲线的交点；通过 \({P}_{1},\ldots ,{P}_{n}\) 的二次曲线线性系统.

\noindent
\textbf{2. 三次曲线与群运算律}\; 曲线 \(\left( {{y}^{2} = x\left( {x - 1}\right) \left( {x - \lambda }\right) }\right)\) 无有理参数化.线性系统 \({S}_{d}\left( {{P}_{1},\ldots ,{P}_{n}}\right)\) ；通过8个“一般位置”点的三次曲线铅笔；三次曲线上的群律；帕斯卡神秘六边形.第一部分附录.曲线及其亏格(genus)

非奇异平面三次曲线在 \(\mathbb{C}\) 上的拓扑；曲线亏格的非正式讨论；拓扑、微分几何、模空间、数论、莫德尔-韦尔-法尔廷斯定理.

\section*{第二部分\quad 仿射簇的范畴}
\noindent
\textbf{3. 仿射簇与零点定理(Nullstellensatz)}\; 诺特环，希尔伯特基定理； \(V\) 与 \(I\) 的对应关系，不可约代数集，Zariski拓扑，零点定理陈述；不可约超曲面.诺特正规化及零点定理证明；归约到超曲面.

\noindent
\textbf{4. 簇上的函数}\; 坐标环与多项式映射；态射与同构；仿射簇.有理函数域与有理映射；占优有理映射及有理映射的复合；标准开集；椭圆曲线上的加法律是态射.

\section*{第三部分\quad 应用}
\noindent
\textbf{5. 射影簇与双有理等价}\; 动机:存在严格大于任何仿射簇的簇；齐次 \(V - I\) 对应；射影与仿射的对比.例子:二次曲面；Veronese曲面.双有理等价，有理簇；每个簇双有理等价于超曲面；积簇.

\noindent
\textbf{6. 切空间与非奇异性，维数}\; 动机:隐函数定理，簇与流形.仿射切空间的定义；非奇异点稠密.切空间与 \(m/{m}^{2}\) ，切空间是内在的； \(X = \operatorname{tr}{\deg }_{k}k\left( X\right)\) 的维数.通过爆破消除奇异点.

\noindent
\textbf{7. 三次曲面上的27条直线}\; 非奇异三次曲面上的直线 \(S\) .通过消元法证明直线的存在；极形式.与给定直线相交的5对直线. \(S\) 是有理的.27条直线的经典构型.黑塞矩阵.所有直线均为有理直线的情况.

\noindent
\textbf{8. 结语}\; 历史与社会学.选题，深奥评论与技术注释.代替前言；致谢与点名.